% TOP_PROBLEM: {"id": 20002546, "title": "Closure-area limit sets and geometric phase on the shape sphere", "category_id": 16, "category": {"id": 16, "name": "physics", "display_name": "Mathematical Physics", "description": "Problems at the intersection of mathematics and physics.", "slug": "physics", "order_index": 16, "created_at": "2026-07-31T15:26:25.671Z"}, "status": "partially_solved", "statusQualification": "The last question concerns realizable dynamical shape curves, not arbitrary loops for which homotopy already provides a deformation. The source leaves the topology and deformation class to be specified. The damaged area notation A0 is repaired. The catalogue title/status was paired with a machine-generated research attempt; that attempt is not peer-reviewed evidence of settlement. The statement covers the original source questions rather than only the title’s subcase."}
% FIRST_POSTED: 2026-10-11
% ENTRY_KIND: statement_only
% SUBMISSION: {"kind": "statement_only", "category": "top_problems", "problem_id": "20002546", "proposal_type": "established", "source": "American Institute of Mathematics, Variational Methods in Celestial Mechanics; notes by Kuo-Chang Chen, version 22 June 2003, source item 17. https://aimath.org/WWN/varcelest/varcelest.pdf", "status_evidence": "Inherited source-catalogue status partially_solved preserved. Background has machine-generated research claims explicitly marked not peer reviewed; these are not independent status evidence. Primary source accessed 11 October 2026; compared against the supplied catalogue excerpt. Historical wording and partial-result remarks were checked for scope; no comprehensive current-literature or proof-verification audit was performed.", "status_checked_on": "2026-10-11", "duplicate_check": "Full 15,773-record source registry checked for identical normalized statements; available site registry checked for identical normalized title/statement and canonical-ID collision. Canonical ID is absent from the site registry. No mathematical-equivalence literature audit performed.", "attribution": "Wu-Yi Hsiang. Statement preparation: OpenAI GPT-6 Codex assistant, generation 11 October 2026; source-matching and notation editing by the assistant, no independent human mathematical verification."}
% SOURCE_STATEMENT: "Problem 17. Problems on the shape space for the three-body problem: (Wu-Yi Hsiang) - Are there solutions whose α-(or ω-)limit set is a limit cycle? - The area of the unit shape sphere is π. Conjecture: for any 0 < A 0 ≤ π, there is a solution whose shape curve has its closure with area A0.- Can two closed shape curves with the same homotopy class be deformed to one another by closed shape curves? Conjecture: No. But then determine connected components."
% Preparation and final self-audit: CONTRIBUTING.md prompts C and D; source/audit artifacts retained in the local physics_statements_2026_10_11 batch.
% Imported from: physics_statements.tex; UnsolvedMath ID 20002546
% !TeX program = lualatex
\documentclass[11pt]{article}
\usepackage[margin=25mm]{geometry}
\usepackage{fontspec}
\setmainfont{Latin Modern Roman}
\usepackage{amsmath,amssymb,mathtools}
\usepackage{unicode-math}
\setmathfont{Latin Modern Math}
\usepackage{xurl,hyperref}
\hypersetup{colorlinks=true,urlcolor=blue}
\setcounter{secnumdepth}{0}
\setlength{\parindent}{0pt}
\setlength{\parskip}{5pt}
\setlength{\emergencystretch}{3em}
\newcommand{\sref}[2]{\hyperlink{source:#1:#2}{[S#2]}}
\newcommand{\cataloguescope}[1]{\paragraph{Recorded target.}#1}
\newcommand{\R}{\mathbb R}
\newcommand{\C}{\mathbb C}
\newcommand{\N}{\mathbb N}
\newcommand{\Z}{\mathbb Z}
\newcommand{\Q}{\mathbb Q}
\newcommand{\F}{\mathbb F}
\newcommand{\eps}{\varepsilon}
\begin{document}
% UnsolvedMath ID 20002546; source catalogue code AIM-PHYSICS-0054
\setcounter{section}{20002545}
\section{Closure-area limit sets and geometric phase on the shape sphere}\label{20002546}
{\small UnsolvedMath ID 20002546 · Mathematical Physics}
\subsection{Definitions and mathematical statement}

\paragraph{Shared notation.}$\mathbb N=\{1,2,\ldots\}$; the standard number fields have their usual meanings. All additional conventions are stated locally.


\paragraph{Definitions and notation.}In the Newtonian $n$-body problem, positions $q_i\in\mathbb R^d$ and masses $m_i>0$ obey $m_i\ddot q_i=\nabla_{q_i}U$, where $U(q)=\sum_{i<j}m_im_j/|q_i-q_j|$ and the gravitational constant is normalized to one. Here $T=\tfrac12\sum_i m_i|\dot q_i|^2$, $E=T-U$, and a collision has $q_i=q_j$. The center of mass can be fixed at zero; moment of inertia is $I=\sum_i m_i|q_i|^2$. For the planar three-body problem, a shape is a configuration modulo translation, rotation and positive scaling. Its normalized shape space is a sphere, with collisions marked and collinear configurations on an equator. A syzygy is a collinear configuration; its symbol records which body is in the middle. For a trajectory, its $\omega$-limit set consists of accumulation points as forward time tends to its endpoint at infinity; its $\alpha$-limit set is defined backward in time. A limit cycle is a closed dynamical orbit. A shape curve is the image of a motion in the shape sphere; its closure has area measured with the source’s sphere normalization, whose total area is $\pi$.

\paragraph{Catalogue wording retained for traceability.}
For the three-body shape space: determine whether a solution can have an $\alpha$- or $\omega$-limit set that is a limit cycle. Prove or refute that for every $0<A_0\le\pi$ there is a solution whose shape-curve closure has area $A_0$. Determine whether two closed shape curves in the same homotopy class can be joined through closed shape curves of solutions; the source conjectures no, and asks for the connected components.
\paragraph{Scope and source qualification.}The last question concerns realizable dynamical shape curves, not arbitrary loops for which homotopy already provides a deformation. The source leaves the topology and deformation class to be specified. The damaged area notation A0 is repaired. The catalogue title/status was paired with a machine-generated research attempt; that attempt is not peer-reviewed evidence of settlement. The statement covers the original source questions rather than only the title’s subcase.
\subsection{Short English statement}
Study shape-space limit cycles, possible closure areas and connected components of closed three-body shape curves.
\subsection{Sources}
\begin{itemize}
\item[\hypertarget{source:20002546:1}{S1}] ULAM.AI and the UnsolvedMath Contributors, UnsolvedMath, record 20002546 (AIM-PHYSICS-0054). Catalogue curation under CC BY 4.0; definitions and English exposition added locally.
\url{https://huggingface.co/datasets/ulamai/UnsolvedMath}.
\item[\hypertarget{source:20002546:2}{S2}] American Institute of Mathematics, Variational Methods in Celestial Mechanics; notes by Kuo-Chang Chen, version 22 June 2003, source item 17.
\url{https://aimath.org/WWN/varcelest/varcelest.pdf}.
\end{itemize}
\label{end:20002546}
\end{document}
